\section{Durchführung}
\label{sec:Durchführung}

\subsection{Aufbau}
Dieser Versuch besteht aus einem Funktionsgenerator, welcher ein Eingangssignal einspeist, einem Oszilloskop zur Darstellung des Signals, einem RLC-Meßgerät zur Messung der verschiedenen Meßgrößen und verschiedenen Koaxialkabeln mit unterschiedlicher Länge und und unterschiedlichem Widerstand. 
Für die erste Messung wird der Funktionsgenerator mit einem kurzen Kabel an das Oszilloskop und einem weiteren Koaxialkabel mit offenem Ende angeschlossen. Hierbei wird ein Eingangssignal vom Funktionsgenerator eingespeist und auf dem Oszilloskop werden die Peaks zur Bestimmung der Länge untersucht. Dabei kann auch die Dämpfungskonstante $\alpha$ bestimmt werden.
In der zweiten Messung werden die Leitungskonstanten $R, L, C$  in Abhängigkeit der Frequenz bestimmt, außerdem wird der Winkel der Phasenverschiebung $\Phi$ aufgenommen mithilfe des RLC-Meßgeräts. Durch geschicktes in Reihe oder parallel schalten kann dann die jeweilige Meßgröße aufgenmommen werden. 
In der dritten Messreihe wird das Koaxialkabel wieder an das Oszilloskop angeschlossen und es werden verschiedene nicht bekannte Abschlusswiderstände an das Kabel angeschlossen. Durch den Verlauf der Signalspannung am Oszilloskop sollen die Abschlusswiderstände bestimmt werden. Hierbei werden 5 verschiedene Abschlusswiderstände verwendet. Relative Dielektrizitätskonstante?
In der letzten Messung werden dann zwei Koaxialkabel in Reihe geschlossen, hierbei werden einmal zwei $50\,\Omega$ Kabel verwendet und einmal ein $50\,\Omega$ und ein $75\,\Omega$ Kabel verwendet, mithilfe der Peaks am Oszilloskop sollen dann die Reflexionsfaktoren bestimmt werden.

\clearpage